% Gerado por gera_latex.py -- não editar à mão.
\setlength{\absparsep}{18pt}
\begin{resumo}
O Diagnóstico da Primeira Infância é um produto desenvolvido pela Coordenadoria de Pesquisa, Avaliação e Política Pública do Instituto Pereira Passos em parceria com a Casa Civil. Ele serve de apoio ao desenho, implementação e monitoramento da nova Política Integrada da Primeira Infância. Enquanto um produto vivo, ele pode e deve ser atualizado conforme novos dados sejam disponibilizados ou novas necessidades surjam. Esta versão em relatório reproduz o conteúdo do website oficial, com alterações de organização e visual apenas.

A maior tendência é a clara redução na participação das crianças até 72 meses (0 a 5 anos) na composição da população. Esse cenário parece primariamente relacionado à significativa redução de nascimentos, tendência que deve continuar a impactar indicadores nos próximos anos. Os dados também apontam melhorias em áreas relativas à prevenção de mortes evitáveis, uma manutenção de padrões de desigualdade social significativos entre as regiões da cidade e a persistência de um alto percentual da população alvo em territórios com altos índices de violência ou em condições de privação socioeconômica.

Ressaltamos aqui que a produção de dados estatísticos e análises ao nível intramunicipal sempre é um desafio e impõe limites à análise. A maior parte da produção estatística brasileira tem como menor unidade de análise o município e as grandes exceções, como os Censos Demográficos, são esporádicos e muitas vezes não atendem as necessidades totais da gestão pública. Por essa razão, com frequência os dados foram mobilizados de fontes distintas, não sendo concebidos originalmente para comparação, ou derivam de registros administrativos que não se destinavam prioritariamente para a análise.

Por fim, os processos sociais aqui apontados podem atuar tanto como causa como consequência de outras variáveis – de ordem econômica, tecnológica, demográfica, etc. Além disso, o período coberto abrange duas significativas crises de saúde pública – a emergência da Zika em 2015-2016 e a Covid-19 em 2020-2021 – além de um período de recessão econômica entre 2015-2016, todas variáveis explicativas significativas, para além do contexto particular da cidade do Rio de Janeiro. Assim, a maior recomendação é que os dados levantados sirvam como um mapa que possa orientar a construção de indicadores socialmente sensíveis aos diversos contextos e não como metas em si mesmas.

\noindent\textbf{Palavras-chave}: Primeira infância. Indicadores sociais. Políticas públicas. Rio de Janeiro (RJ).
\end{resumo}
